\subsection*{Part B: precise scope of the remaining classification demand}
B.009 is an exact necessary-and-sufficient tensor criterion for arbitrary
$R\subset S$, separating injectivity from existence. B.013 completely
classifies every finite base ring with arbitrary extension; B.014
classifies every field base and reduces torsion-free extensions of
infinite domains to rational coefficient overrings. B.017 classifies
all base rings for which every extension works: exactly those with no
finite maximal residue field. B.019 proves the finite dual-basis
criterion, in particular the projective case; B.020 covers every
Noetherian normal domain with every torsion-free coefficient extension.

B.021--B.026 prove projectivity for every finite-conductor subring of
a finite-residue discrete valuation ring, including arbitrary residue
extensions, mixed characteristic, incomplete rings and conductor gaps.
They include the two formerly conjectured nonflat cusp examples
recorded by Elliott and refute those particular nonflatness claims.
B.027 proves the required polynomial localization directly. B.030--B.032
then give a complete structural classification for every one-dimensional
Noetherian domain with finite normalization and every torsion-free
coefficient extension. At each finite-residue maximal ideal, the rational
coefficient overring must equal the base local ring, or equal its fraction
field when the normalization has exactly one branch. Finite intermediate
extensions fail tensor injectivity; other extensions at a multibranch
point fail even multiplicative existence. All implications are proved.

B.033 extends polynomial localization to every Noetherian dimension.
B.034 identifies the punctured-spectrum coefficient ring and its depth
test. B.035--B.036 complete the classification in every dimension with
finite normalization: at higher-dimensional finite-residue local rings,
the rational coefficient overring must intersect the punctured-spectrum
ring in the base ring itself. For fraction-field coefficients this is
exactly depth at least two. B.037 gives an explicit two-dimensional
nonzero tensor kernel, with a full finite certificate.

B.033--B.036 extend the complete classification to every Noetherian
domain with finite normalization, in all dimensions and for all
torsion-free coefficient extensions. At a finite-residue maximal
ideal of height at least two, the criterion is
$L_{\mathfrak m}\cap R_{\mathfrak m}^{\sharp}=R_{\mathfrak m}$,
where $R_{\mathfrak m}^{\sharp}$ is explicitly defined by saturation
by powers of the maximal ideal. For fraction-field coefficients
this is exactly depth at least two. B.038--B.040 also exclude every
finite-residue Akizuki--Reid family with nonfinite normalization,
by an explicit tensor kernel. B.041 proves every infinite valuation
domain, of arbitrary rank, works for every torsion-free extension.
B.042 proves the same for every infinite Krull domain, with no
Noetherian assumption. B.043 completely classifies all coefficient
extensions for the non-Noetherian cusps $\mathbb F_q+t^a l[[t]]$
with $l$ infinite: only rational coefficient ring $L=R$ works.

The broader structural classification for arbitrary specified unital
inclusions is not proved. The remaining domain, torsion-free scope
is outside all these classes: Noetherian domains with nonfinite
normalization other than the proved derivation-obstruction families,
and non-Noetherian domains outside the established positive classes.
B.044--B.046 now completely settle fraction-field coefficients
for every one-dimensional Noetherian domain: each finite-residue
local completion must be a domain, equivalently its normalization
must be finite and local. Thus no one-dimensional fraction-field
case remains undecided. Arbitrary proper rational coefficient
overrings in the nonfinite-normalization case remain a precise
residual class; the Akizuki--Reid examples already excluded above
are not listed as undecided. Zero-divisor bases and extensions with base
torsion have the exact criterion and the finite/universal
classifications, but no general invariant classification here.
The exact tensor criterion is not presented as a classification by
other ring-theoretic invariants of these remaining cases.

A 2016 source initially suggested a stronger classification, but its
official arXiv record states that it was withdrawn on 2025-06-10
because of critical errors; no result from it is invoked. Elliott's
2010 questions describe that paper's date, not current status. The
Jiang et al. (2026) counterexample was opened in both PDF and author
LaTeX, checked, and generalized in B.016 and B.025. Its Section 3.4
contains the explicit two-branch example, not the general positive or
branch-classification theorem proved here. This comparison does not
assert priority or absence from other literature.

No historical priority or independence assertion is made. L6 is
settled as the exact structural comparison in B.012.
